\section{Positive domination by cap measures}\label{p:cap-section}

This section proves the sample bound by propagating the joint law of the
signal and each memory state forward through the stream.  We bound each
such law by a positive sum of uniform measures on spherical caps.  A term
is charged for its mass and for a power of the reciprocal of its radius.
The main step shows that a block of observations increases the total charge
by a controlled factor.  The passage from bounds on test functions to an
actual dominating measure requires compactness in the norm of $L^1$.

We prove Theorem~\ref{p:regression} by a forward argument in the
finite-state experiment of Section~\ref{p:geometry-model}. Use that
section's fixed-seed convention and Borel kernels from
Lemma~\ref{p:borel-kernels}. All bounds below are uniform in the fixed
rules; their scalar probability conclusions are averaged over the shared
seed at the end.

\subsection{The spherical experiment and the projection input}

Write $\sigma$ for uniform probability on $S^{d-1}$, and set
\begin{equation}\label{p:cap-parameters}
 n=d-1,\qquad k=p=\lfloor d/16\rfloor,\qquad
 \alpha=\lambda=\frac n2,\qquad \theta=1-\frac1p.
\end{equation}
We take $d$ sufficiently large that $p\ge2$ and $k,p\ge d/32$.
For $c\in S^{d-1}$ and $0<r\le2$, define
\[
 B(c,r)=\{s\in S^{d-1}:\norm{s-c}\le r\},\qquad
 \mu_{c,r}(ds)=\frac{\1_{B(c,r)}(s)}{\sigma(B(c,r))}\,\sigma(ds).
\]
These are spherical caps and their normalized probability measures.
For an ambient ball in $\R^d$ we use the distinct notation
$D(b,t)=\{u\in\R^d:\norm{u-b}\le t\}$.

Lemma~\ref{p:cap-mass} supplies the two-sided estimate
\begin{equation}\label{p:cap-volume-bound}
 (ar)^n\le\sigma(B(c,r))\le(A_0r)^n,
 \qquad a=(2\pi)^{-1},\quad A_0=1.
\end{equation}
In particular every $\mu_{c,r}$ is well defined.

Let $X$ be a $k\times d$ matrix with independent standard Gaussian
entries and put
\[
 H=(XX^{\mathsf T})^{1/2},\qquad P=H^{-1}X
\]
on the probability-one event that $X$ has full row rank.  Define $P$
arbitrarily on the remaining null set.  We may make this definition
measurable.  On the full-rank event $PP^{\mathsf T}=I_k$ and $X=HP$;
the law of $P$ is
invariant under right orthogonal multiplication.  Let
$v_k$ be the Lebesgue volume of the unit ball of $\R^k$ and let
$m_k(dz)=v_k^{-1}\1_{\{\norm z\le1\}}\,dz$ be uniform probability on
that ball.

The use of independent replicas for high projection moments is related
to the argument of Sharan, Sidford, and Valiant
\cite[Section~7, Lemma~12]{SSV2019}. We use the normalized moment theorem
(Theorem~\ref{p:frame-positive}) with the parameters in
\eqref{p:cap-parameters}. Its inverse-distance margin satisfies
$\lambda-(p-2)-k\ge3d/8$. Below we verify its unit support, mass,
local-growth, independence, and absolute-continuity hypotheses for each
finite positive test measure before invoking its normalized density bound.

\subsection{One block: from bounded tests to a dominating measure}

Fix a parent cap $B=B(c,r)$.  Let $\mathcal W$ be a finite set and let
$K_w(X,Y)\in[0,1]$, for $w\in\mathcal W$, be Borel functions of a
$k\times d$ matrix and a vector in $\R^k$ with $\sum_wK_w\le1$.
Define
\begin{equation}\label{p:cap-reference-probabilities}
 \begin{split}
 f_w(s)&=\E_X K_w(X,Xs),\\
 q_w&=\E_X\int K_w(X,Xc+rHz)\,dm_k(z).
 \end{split}
\end{equation}
The functions $f_w$ are measurable.  The variables $z$ in the second
line are independent of $X$ and are used only to define a reference law.
The kernel still receives the full matrix $X$ and the radial factor $H$
through the displayed label.  In particular, $\sum_wq_w\le1$.

\begin{proposition}[Positive domination after one block]
\label{p:cap-domination}
For each $w\in\mathcal W$, there are countably many caps
$B(c_i,\rho_i)$ with $0<\rho_i\le r$ and coefficients $a_i\ge0$ such
that
\begin{equation}\label{p:cap-domination-bound}
 f_w\mu_{c,r}\le\sum_i a_i\mu_{c_i,\rho_i},
 \qquad
 \sum_i a_i(r/\rho_i)^\alpha\le\exp(Cd)q_w^\theta.
\end{equation}
The first inequality is an inequality of finite positive measures on
$S^{d-1}$: it holds after evaluating both sides on every measurable set.
The constant $C$ is independent of the parent cap and of the kernels.
\end{proposition}

\begin{proof}
We prove a bound for nonnegative tests and then convert it into a
countable positive domination.

\smallskip\noindent
\emph{Cap tests and their local mass.}
Put $\rho_j=r2^{-j}$ and choose a finite $\rho_j/2$-net $N_j$ of the
sphere for every $j\ge0$. Take as a dictionary the parent cap, indexed
by $*$, together with all $B(c',\rho_j)$, $c'\in N_j$. Repetitions may
have separate indices. For this countable index set $I$, write
$B_i=B(c_i,\rho_i)$ and
\[
 \phi_i=\frac{\1_{B(c_i,\rho_i)}}{\sigma(B(c_i,\rho_i))},
 \qquad b_i=(r/\rho_i)^\alpha.
\]
Thus $\phi_*\,d\sigma=\mu_{c,r}$, every $\phi_i$ is a probability
density, $b_*=1$, and the weights at level $j$ are $2^{\alpha j}$.
If $h\in L^\infty(\sigma)$ is nonnegative and satisfies
\begin{equation}\label{p:cap-dual-constraints}
 \int h\phi_i\,d\sigma\le b_i\qquad(i\in I),
\end{equation}
then its integral against any positive dictionary sum is at most that
sum's weighted cost. We will bound $\int h f_w\,d\mu_{c,r}$ for every
such test.

Let $\nu$ be the image of $h\mu_{c,r}$ under $s\mapsto(s-c)/r$.
It has unit-ball support and mass at most one by the parent constraint.
We claim that, for an absolute $D\ge1$,
\begin{equation}\label{p:cap-test-ball-bound}
 \nu(D(b,t))\le D^d t^\lambda
 \qquad(b\in\R^d,\ t>0).
\end{equation}
For $0<t\le1/4$, suppose the ball meets the rescaled cap, and choose a
point $s_0\in B$ whose image lies in it. The corresponding sphere points
are all within $2rt$ of $s_0$. Set
\[
 j=\left\lfloor\log_2\frac1{4t}\right\rfloor,\qquad \rho=\rho_j;
 \qquad 4rt\le\rho<8rt,\quad \rho\le r.
\]
A center in $N_j$ within $\rho/2$ of $s_0$ gives a dictionary cap
containing those points, since $2rt+\rho/2\le\rho$. Its constraint and
the two-sided cap estimate imply
\begin{align*}
 \nu(D(b,t))
 &\le\frac{\sigma(B_i)}{\sigma(B)}\int h\,d\mu_{c_i,\rho_i}\\
 &\le(A_0/a)^n(\rho/r)^{n-\alpha}
 \le(A_0/a)^n(8t)^\lambda.
\end{align*}
An empty intersection has zero mass. For
$t>1/4$, use the mass bound and $1\le4^\lambda t^\lambda$.
These estimates prove \eqref{p:cap-test-ball-bound}.

\smallskip\noindent
\emph{The quantitative test estimate.}
All hypotheses of Theorem~\ref{p:frame-positive} now hold for $\nu$.
Indeed, the test, parent cap, and measure are fixed before the fresh
matrix $X$ is drawn. For every orthonormal $k$-frame $P$, the coordinate
density in Lemma~\ref{p:coordinate-density} gives $P_\#\sigma\ll dz$.
Because $h\mu_{c,r}\ll\sigma$, its projection is also absolutely
continuous, and the affine rescaling
\[
 P((s-c)/r)=(Ps-Pc)/r
\]
gives the same conclusion for $P_\#\nu$. The theorem therefore supplies
a jointly measurable density satisfying, for almost every $X$,
\[
 P_\#\nu=g_Xm_k,\qquad g_X(z)=0\text{ for }\norm z>1,
 \qquad \E_X\int g_X^p\,dm_k\le\exp(C_Ddp).
\]
The density is relative to $dm_k=dz/v_k$ on the unit ball; $D$, and
hence the new constant, is absolute. Since $s=c+ru$ gives
$Xs=Xc+rHPu$, Tonelli's theorem and H\"older's inequality on the product
law of $X$ and $z$ yield
\begin{equation}\label{p:cap-dual-test}
\begin{split}
 \int h f_w\,d\mu_{c,r}
 &=\E_X\int K_w(X,Xc+rHz)g_X(z)\,dm_k(z)\\
 &\le\left(\E_X\int g_X^p\,dm_k\right)^{1/p}
 \left(\E_X\int K_w(X,Xc+rHz)^{p'}\,dm_k\right)^{1/p'}\\
 &\le\exp(Cd)q_w^\theta,\qquad p'=p/(p-1).
\end{split}
\end{equation}
The last line uses $0\le K_w\le1$. Thus every test satisfying
\eqref{p:cap-dual-constraints} obeys \eqref{p:cap-dual-test}.

\smallskip\noindent
\emph{Representation by positive cap sums.}
The geometric estimate is complete; the remaining step uses only the
dictionary and its test bound. If $q_w=0$, the feasible test $h=1$ shows
that $f_w\mu_{c,r}=0$, and the empty sum suffices. Otherwise set
$K=\exp(Cd)q_w^\theta>0$, with the constant from
\eqref{p:cap-dual-test}, and define in real $L^1(\sigma)$
\[
 \mathcal C_K=\left\{\sum_{i\in I}a_i\phi_i:
 a_i\ge0,\ \sum_i a_ib_i\le K\right\}.
\]
Each series converges in $L^1$, since $b_i\ge1$ and
$\norm{\phi_i}_{L^1}=1$. This convex set is norm compact. To see this,
take a sequence and choose a coefficient array for each term. A diagonal
subsequence converges coefficientwise, since $0\le a_i\le K/b_i$.
The limiting array has cost at most $K$ by Fatou's lemma. There are
finitely many indices through every level $J$, and uniformly over all
admissible arrays,
\begin{equation}\label{p:cap-compact-tail}
 \left\|\sum_{j>J}\sum_{i\text{ at level }j}a_i\phi_i\right\|_{L^1}
 =\sum_{j>J}\sum_{i\text{ at level }j}a_i
 \le K2^{-\alpha(J+1)}.
\end{equation}
The same estimate holds for the limiting array. Finite-level convergence
followed by \eqref{p:cap-compact-tail} gives convergence in $L^1$.
This proves sequential compactness, and hence compactness in the metric
space $L^1(\sigma)$
\cite[Section~1.3, Theorem~1.37, p.~29]{SchnaubeltFA2025}.

Its signed downward closure
\[
 \mathcal D_K=\{u\in L^1(\sigma):u\le v\text{ a.e. for some }v\in\mathcal C_K\}
\]
is convex, contains every nonpositive $L^1$ function, and is norm closed.
In fact, if $u_m\to u$ and $u_m\le v_m\in\mathcal C_K$, compactness
provides a subsequence $v_m\to v\in\mathcal C_K$; then
\[
 \norm{(u-v)_+}_{L^1}
 \le\norm{u-u_m}_{L^1}+\norm{v_m-v}_{L^1}\longrightarrow0.
\]
Put $F=f_w\phi_*$. If $F\notin\mathcal D_K$, strict separation of the
closed convex set $\mathcal D_K$ and the compact singleton $\{F\}$
\cite[Section~5.2, Theorem~5.20(b), pp.~98--99]{SchnaubeltFA2025}
gives a continuous real functional $\ell$ with
\[
 \ell(F)>\sup_{u\in\mathcal D_K}\ell(u).
\]
A probability measure is localizable
\cite[Theorem~211L]{FremlinTaxonomy}, so real $L^1$--$L^\infty$
duality~\cite[Theorem~243G(b)]{FremlinFunctionSpaces} represents
$\ell(u)=\int hu\,d\sigma$ for a real $h\in L^\infty(\sigma)$.
The finite supremum forces $h\ge0$: if $h\le-\delta$ on a set $E$
of positive measure, the functions $-t\1_E\in\mathcal D_K$ would make
it infinite. Define
\[
 S_h=\sup_i\frac{\int h\phi_i\,d\sigma}{b_i},\qquad
 0\le S_h\le\norm h_{L^\infty}.
\]
Then
\begin{equation}\label{p:cap-support-functional}
 \sup_{u\in\mathcal D_K}\int hu\,d\sigma=KS_h.
\end{equation}
The cost bound gives the upper inequality, and the single terms
$(K/b_i)\phi_i$ give the reverse inequality by taking the supremum over
$i$. This does not require the supremum to be attained. If $S_h=0$,
the parent term gives $\int h\phi_*\,d\sigma=0$, and $0\le f_w\le1$
implies $\ell(F)=0$, contrary to strict separation. If $S_h>0$, the
function $h/S_h$ satisfies \eqref{p:cap-dual-constraints}, so
\eqref{p:cap-dual-test} again contradicts separation. Therefore
$F\in\mathcal D_K$. An admissible series dominates its density almost
everywhere; integrating gives the asserted measure inequality and cost
bound. Its total mass is at most $K$.
\end{proof}

\subsection{Propagation, bounded stopping, and the accuracy endpoint}

The preceding proposition applies to each cap in a joint-state bound.
The positivity of all measures makes the iteration linear before any
estimate is taken.

\begin{lemma}[Cost of a forward block]\label{p:cap-forward-step}
Let $\mathcal V$ and $\mathcal W$ be finite state sets with
$\abs{\mathcal W}\le2^M$.  For each $v\in\mathcal V$, let
$K_{v,w}$ be block kernels of the type in
\eqref{p:cap-reference-probabilities}, with likelihoods $f_{v,w}$.
Let $\eta_v$ be finite positive measures on the sphere and suppose
\[
 \eta_v\le\sum_i a_{v,i}\mu_{c_{v,i},r_{v,i}},\qquad
 \mathcal P=\sum_{v,i}a_{v,i}(2/r_{v,i})^\alpha<\infty,
\]
where each sum is countable, $a_{v,i}\ge0$, and $0<r_{v,i}\le2$.
Define the outgoing measures by
\[
 \eta'_w(E)=\sum_v\int_E f_{v,w}(s)\,\eta_v(ds).
\]
Then the $\eta'_w$ admit countable positive cap dominations whose total
cost $\mathcal P'$ satisfies
\begin{equation}\label{p:cap-cost-multiplier}
 \mathcal P'\le\mathcal P\exp(Cd)2^{M/p}.
\end{equation}
\end{lemma}

\begin{proof}
Multiplication by a nonnegative measurable function preserves a measure
inequality, as follows first for simple functions and then by monotone
convergence.  Summing over incoming states and cap terms therefore gives
\[
 \eta'_w\le\sum_{v,i}a_{v,i}f_{v,w}\mu_{c_{v,i},r_{v,i}}.
\]
Apply Proposition~\ref{p:cap-domination} separately to every term.
For a fixed incoming cap, write $q_{v,i,w}$ for its reference
probabilities.  They depend on that cap as well as on the starting state,
and satisfy $\sum_wq_{v,i,w}\le1$.  H\"older's inequality for the finite
sum gives
\[
 \sum_wq_{v,i,w}^{\theta}
 \le\abs{\mathcal W}^{1/p}\left(\sum_wq_{v,i,w}\right)^\theta
 \le2^{M/p}.
\]
If the outgoing terms for this parent have coefficients $\beta_j$ and radii
$\rho_j$, their contribution to the global cost is bounded by
\begin{align*}
 \sum_j a_{v,i}\beta_j(2/\rho_j)^\alpha
 &=a_{v,i}(2/r_{v,i})^\alpha
       \sum_j \beta_j(r_{v,i}/\rho_j)^\alpha\\
 &\le a_{v,i}(2/r_{v,i})^\alpha\exp(Cd)q_{v,i,w}^\theta.
\end{align*}
Sum first over $w$ and then over the countably many incoming terms to
obtain \eqref{p:cap-cost-multiplier}.  Tonelli's theorem applies
throughout because every summand is nonnegative.

At each of finitely many block boundaries there are only countably many
state--cap terms.  Choosing a dominating sum for each produces a
deterministic proof representation, with another countable sum at the
next boundary.  There is no choice indexed by the random signal or by
the fresh matrix.  The state measures here remain unnormalized joint
measures; the argument never divides them by their total masses.
\end{proof}

\begin{lemma}[Terminal success under a cap]\label{p:cap-terminal}
Let $Q$ be any probability law on unit-vector outputs, independent of
the signal, and put
$h_Q(s)=Q\{z:\arccos\inner{z}{s}\le\epsilon\}$ for
$0<\epsilon\le1/10$.  For every cap $B(c,r)$,
\begin{equation}\label{p:cap-terminal-bound}
 \int h_Q\,d\mu_{c,r}\le\exp(Cd)(\epsilon/r)^\alpha.
\end{equation}
\end{lemma}

\begin{proof}
For a fixed unit output $z$, write $\vartheta=\arccos\inner{z}{s}$.
Angular success implies $\norm{s-z}=2\sin(\vartheta/2)\le\epsilon$.  If
$\epsilon<r$, Lemma~\ref{p:cap-mass} bounds its probability under
$\mu_{c,r}$ by
\[
 \frac{\sigma(B(z,\epsilon))}{\sigma(B(c,r))}
 \le(A_0/a)^n(\epsilon/r)^n
 \le\exp(Cd)(\epsilon/r)^\alpha.
\]
If $\epsilon\ge r$, use probability at most one and
$(\epsilon/r)^\alpha\ge1$.  Integration over $Q$ proves the result.
\end{proof}

\begin{proposition}[A fixed stopping index]\label{p:cap-fixed-stop}
For all sufficiently large $d$, let a learner obey the observation and
information restrictions in Section~\ref{p:geometry-model}, with
$0<\epsilon\le1/10$, any nonnegative integer $M$ bits of state, and horizon $T$,
and let $\mathsf T$ be its stopping index.  For every fixed integer
$0\le t\le T$,
\begin{equation}\label{p:cap-fixed-stop-explicit}
 \Prob\{\mathsf T=t\text{ and }\arccos\inner{\widehat S}{S}\le\epsilon\}
 \le\exp(Cd)(\epsilon/2)^\alpha
       \left[\exp(Cd)2^{M/p}\right]^{\lceil t/k\rceil}.
\end{equation}
In particular, for a sequence $M(d)=o(d^2)$ the right side is eventually
at most
\begin{equation}\label{p:cap-fixed-stop-asymptotic}
 \exp\bigl(Cd(1+\lceil t/k\rceil)\bigr)(\epsilon/2)^\alpha,
\end{equation}
with an absolute $C$ and a dimension threshold that may depend on the
memory sequence.
\end{proposition}

\begin{proof}
Fix $t>0$ throughout the argument and split its first $t$ indices into
$b=\lceil t/k\rceil$ consecutive blocks of length at most $k$.  At a
prefinal boundary retain only paths that have not stopped and are eligible
to continue.  In the final block retain exactly the paths stopping at
$t$, with their terminal state.  For a fixed starting state and a supplied
block of pairs, integrating the internal randomness gives the probability
of each retained ending state.  These are Borel kernels with sum at most
one, since all omitted paths are discarded.  A shorter block is padded
by independent unused Gaussian rows, which the kernels ignore.  Its
likelihood for a signal $s$ is therefore $f_{v,w}(s)$ of
\eqref{p:cap-reference-probabilities}. Conditional on the fixed seed,
the fresh Gaussian feature rows in the block are jointly independent of the signal and entering state,
indeed of the signal, initialization, past rows, and past coins. Given the
signal, the past affects this likelihood only through the starting state
and the fixed index.

At any counted boundary, let $\eta_v(E)$ be the probability that
$S\in E$ and the path is counted in state $v$.  These are unnormalized
joint measures and their outgoing laws have exactly the formula in
Lemma~\ref{p:cap-forward-step}. Initially, $B(c,2)=S^{d-1}$ and
$\mu_{c,2}=\sigma$. Write $V_0$ for the initial finite state. For the
fixed shared seed $\xi$, joint initialization
independence gives
\[
 \eta_v^\xi(E)=\pi_\xi(v)\sigma(E),
 \qquad \pi_\xi(v)=\Prob\{V_0=v\mid\xi\},\qquad
 \sum_v\pi_\xi(v)=1.
\]
Thus the initial total cost is one before any discarded paths are removed,
and at most one afterward.
After $b$ applications of that lemma, the counted terminal measures
admit positive cap dominations of total cost at most
$[\exp(Cd)2^{M/p}]^b$.

At each fixed terminal state and index $t$, the output is drawn from a
fixed law using fresh randomness. Apply Lemma~\ref{p:cap-terminal} to each
term of its dominating cap sum.  A term of coefficient $a$ and
radius $r$ contributes at most
$\exp(Cd)(\epsilon/2)^\alpha a(2/r)^\alpha$.
Summing proves the conditional version of
\eqref{p:cap-fixed-stop-explicit}. At $t=0$, use the
initial domination and the same terminal estimate directly.  The index
has been fixed throughout, so only the at most $2^M$ states at each
ending boundary were counted.  No additional state is charged for the
value of $t$.

The countable dominating sums are chosen only after fixing the shared
seed. The bound is uniform in that seed, so the common convention of
Section~\ref{p:geometry-model} allows us to average the original
conditional success probabilities. No measurable choice of cap
decompositions across shared seeds is required.
Finally, for $M(d)=o(d^2)$ we eventually have $M(d)\le d^2$.  Since
$p\ge d/32$, the logarithm of the extra factor $2^{M/p}$ is then at
most $32(\log2)d$.  Absorbing it into the absolute exponential constant
proves \eqref{p:cap-fixed-stop-asymptotic}.
\end{proof}

We now sum the fixed-index estimate over all possible stopping times
and use the common full-data bound to absorb the resulting lower-order
terms.

\begin{proof}[Proof of Theorem~\ref{p:regression} by positive cap domination]
Put $L=\log(1/\epsilon)$.

For the given sequence $M(d)=o(d^2)$, use
\eqref{p:cap-fixed-stop-asymptotic} in all sufficiently large dimensions.
The stopping events are disjoint. Since $k\ge d/32$, $\alpha\ge d/3$,
and $\lceil t/k\rceil\le t/k+1$, summing over $0\le t\le T$ yields
\begin{equation}\label{p:cap-total-success}
 \Prob\{\text{success}\}
 \le(T+1)\exp(C_1d+C_2T-dL/3).
\end{equation}
Here $(\epsilon/2)^\alpha\le\exp(-dL/3)$, and $C_1,C_2$ are absolute.
If success is at least $2/3$, taking logarithms gives
\[
 \frac{dL}{3}
 \le C_1d+C_2T+\log(T+1)+\log(3/2)
 \le C_1d+(C_2+1)T+\log(3/2),
\]
since $\log(T+1)\le T$ for every $T\ge0$.
Lemma~\ref{p:haar-linear} also gives $T>d/8$, and therefore $T\ge1$.
Consequently
\[
 dL\le3(8C_1+C_2+2)T.
\]
This proves the claimed bound with an absolute constant, uniformly over
the full accuracy range. Only the dimension at which the memory factor
is absorbed may depend on the sequence $M(d)$.
\end{proof}
